\Article{368}
Congruentiam aliquam $S\equiv0$ radices seu generalius divisores \emph{aequales} habere dicimus, si per functionis alicuius potestatem dividi possit.

\SourceLeaf{246}
Num congruentia proposita divisores aequales habeat, eodem modo diiudicatur, uti in aequationum theoria. Ponamus
\[
X\equiv\xi^mP
\]
patet fore
\[
\frac{\mathrm{d}X}{\mathrm{d}x}\equiv\xi^{m-1}\left(mP\frac{\mathrm{d}\xi}{\mathrm{d}x}+\xi\frac{\mathrm{d}P}{\mathrm{d}x}\right)
\]
quare $\frac{\mathrm{d}X}{\mathrm{d}x}$ per $\xi^{m-1}$ dividetur. Generaliter sit
\[
X\equiv A^aB^bC^c\text{ etc.}
\]
ubi $A,B,C$ etc. denotant functiones primas diversas, erit
\[
\frac{\mathrm{d}X}{\mathrm{d}x}\equiv X\left(\frac{a\,\mathrm{d}A}{A\,\mathrm{d}x}+\frac{b\,\mathrm{d}B}{B\,\mathrm{d}x}+\frac{c\,\mathrm{d}C}{C\,\mathrm{d}x}+\text{etc.}\right)
\]
unde patet, nisi aliquis numerorum $a,b,c$ etc. per modulum dividatur, $\frac{\mathrm{d}X}{\mathrm{d}x}$ per $A^{a-1}B^{b-1}C^{c-1}$ etc. dividi posse, non autem per $A^a,B^b,C^c$ etc. Hinc sequitur

\textsc{Theorema.} \emph{Si functionum $X$ et $\frac{\mathrm{d}X}{\mathrm{d}x}$ divisor communis maximae dimensionis sit $\xi$, omnes factores primos, quos $\xi$ habet, etiam $X$ habebit et quidem quemvis toties $+1$ vice quoties $\xi$, si igitur $X$ et $\frac{\mathrm{d}X}{\mathrm{d}x}$ sint functiones inter se primae, $X$ nullos factores aequales habebit.}

\Article{369}
Exemplum I. Quaeritur an functio
\[
x^5+3x^4-6x^3+3x-4\ldots(X)
\]
secundum modulum $17$ divisores aequales habeat. Erit
\[
\frac{\mathrm{d}X}{\mathrm{d}x}\equiv5x^4-5x^3-xx+3
\]
Hinc invenitur, functiones $X$ et $\frac{\mathrm{d}X}{\mathrm{d}x}$ inter se esse primas, quare $X$ divisores aequales non habet.

Exemplum II. Sit
\[
X\equiv x^5+6x^4-3x^3-4xx+2x-3(\mathrm{mod.}\,13)
\]
erit
\[
\frac{\mathrm{d}X}{\mathrm{d}x}\equiv5x^4-2x^3+4xx+5x+2
\]
maxima vero functionum $X,\frac{\mathrm{d}X}{\mathrm{d}x}$ communis mensura $\equiv5xx+7x+7$ seu 
\SourceLeaf{247}
multiplicata per $8$: $xx+4x+4$; at quum hic divisor sit $\equiv(x+2)^2$, functio $X$ per $(x+2)^3$ dividi poterit quotiensque (qui est $xx+11$) nullum amplius divisorem duplicem involvit.

\Article{370.\ 371}
Si ex \textsection.\textsection.\ praecc. functio $X$ ita est exhibita $A^aB^bC^c$ etc., ita ut $A,B,C$ etc. inter se sint primae et numeri $a,b,c$ etc. inaequales, resolutio etiam ulterius extendi potest. Sit itaque $X$ functio, quae nullos amplius divisores aequales involvit. Supra vidimus, $x^p-x$ esse productum ex omnibus functionibus primis unius dimensionis. Sit $\xi$ divisor communis maximae dimensionis functionum $X$ et $x^p-x$, erit $\xi$ productum ex omnibus divisoribus ipsius $X$ unius dimensionis, et $\frac X\xi$ huiusmodi divisores non amplius habebit. Quodsi autem inveniatur, functiones $X$ et $x^p-x$ esse inter se primas, $X$ nullum divisorem unius dimensionis habebit adeoque congruentia $X\equiv0$ radices reales non habebit. Porro quoniam $x^{pp}-x$ est productum ex omnibus functionibus primis duarum dimensionum uniusque, divisor communis maximae dimensionis functionum $x^{pp}-x$ et $\frac X\xi$, $\xi'$ involvet omnes divisores ipsius $X$, qui sunt duarum dimensionum. Hinc ulterius progrediendo perspicitur, $X$ hoc modo in factores $\xi,\xi',\xi''$ etc. resolvi, qui continent respective omnes divisores unius, duarum, trium etc. dimensionum.

\Article{372}
Si autem productum ex pluribus functionibus primis eiusdem dimensionis datum est, singulae functiones tentando erui debent. Magnam analogiam habet hoc problema cum eo, quod numerorum compositorum factores quaerere iubet. Hic vero iam a priori determinatur, an functio proposita in factores adhuc discerpi possit. Quum et hic factorum omnium possibilium multitudo sit finita, simili subsidio ut supra uti possumus. Sed huic rei inhaerere nolumus, nam calculator exercitatus principia probe assecutus, quando opus est, facile artificia particularia reperiet.

Progredimur ad aliud caput, scilicet ad considerationem congruentiarum, si modulus non est numerus primus, uti hactenus semper supposuimus. Praesertim vero hic ille casus attentione dignus est, ubi modulus est numeri primi potestas, tum per se tum quod ad aliqua dubia removenda (\textsection.\textsection.\ .\ .) necessarius sit.
